\documentclass[aspectratio=169,11pt]{beamer}

\usetheme{Madrid}
\usecolortheme{default}
\usefonttheme{professionalfonts}
\setbeamertemplate{navigation symbols}{}
\setbeamertemplate{footline}[frame number]

\usepackage{amsmath, amssymb, booktabs}

\title{Education, Skills, Aspirations, And Occupational Sorting}
\subtitle{Gender Study --- Week 3}
\author{Teaching deck draft}
\date{}

\begin{document}

\begin{frame}
  \titlepage
\end{frame}

\begin{frame}{Central question}
\begin{itemize}
    \item How do education, skills, aspirations, and early occupational choices generate later gendered labor-market outcomes?
    \item This is labor economics because sorting shapes skill prices, job access, career ladders, wage growth, authority, and job quality.
    \item Week 3 discipline: separate early choices from later adjustment frictions.
\end{itemize}
\end{frame}

\begin{frame}{Week 2 to Week 3 bridge}
\begin{itemize}
    \item Week 2: household allocation, care, fertility, and expected family roles.
    \item Week 3: those expectations can affect schooling, major choice, aspirations, internships, search, and first jobs.
    \item A flexible first track can be an early choice shaped by expected later constraints.
\end{itemize}
\end{frame}

\begin{frame}{Why sorting is not one mechanism}
\small
\begin{tabular}{p{0.24\linewidth} p{0.55\linewidth}}
\toprule
Channel & Object \\
\midrule
Early choices & schooling, major, track, skills, aspirations, role models, information \\
Entry sorting & first occupation, sector, employer, task, ladder, schedule \\
Later adjustment & switching costs, promotion gates, flexibility constraints, persistence \\
Interactions & anticipated future constraints shape early choices; early choices raise later costs \\
\bottomrule
\end{tabular}
\end{frame}

\begin{frame}{Conceptual decomposition}
\[
G_Y =
\Delta^{early}_Y
+ \Delta^{entry}_Y
+ \Delta^{adjust}_Y
+ \Delta^{returns}_Y
+ \Delta^{interactions}_Y
\]
\begin{itemize}
    \item \(Y\): earnings, wages, authority, hours, promotion, job quality, or welfare.
    \item Early: skills, schooling, aspirations, beliefs, information.
    \item Adjustment: persistence, switching, family timing, mobility, promotion.
    \item Decomposition is not identification; it is a map of claims.
\end{itemize}
\end{frame}

\begin{frame}{Early-choice object}
\[
o_i^0 =
\arg\max_{o \in \mathcal{O}}
\left\{
\mathbb{E}_{i0}\left[\sum_{t=0}^{T}\beta^t w_{ot}(H_{it})\right]
+ A_{io}
- C_{io}
- D^e_{io}
- \kappa^e_{io}
\right\}
\]
\begin{itemize}
    \item \(A_{io}\): aspirations, identity fit, perceived belonging.
    \item \(D^e_{io}\): expected discrimination or exclusion.
    \item \(\kappa^e_{io}\): anticipated future adjustment costs.
\end{itemize}
\end{frame}

\begin{frame}{Early choices channel}
\begin{itemize}
    \item Field of study, training, courses, internships, and first occupation.
    \item Role models and mentors shift information and perceived belonging.
    \item Teacher signals and stereotypes shape perceived comparative advantage.
    \item Competition, risk, and aspirations affect entry into high-variance tracks.
    \item Anticipated discrimination changes choices before employer treatment is observed.
\end{itemize}
\end{frame}

\begin{frame}{Aspirations, competition, information, role models}
\small
\begin{tabular}{p{0.30\linewidth} p{0.48\linewidth}}
\toprule
Mechanism & Observed margin \\
\midrule
Female role-model exposure & economics major choice \\
Competitiveness measure & later track or career choice \\
Teacher feedback / bias & course taking, achievement, field choice \\
Information about returns or fit & applications, major choice, intended field \\
\bottomrule
\end{tabular}
\vspace{0.3cm}
\begin{itemize}
    \item Name the identifying variation before interpreting the result.
\end{itemize}
\end{frame}

\begin{frame}{Anticipated discrimination and field choice}
\begin{itemize}
    \item Expected treatment can enter the supply-side choice problem.
    \item A student may avoid a field because she expects worse hiring, mentoring, grading, promotion, or workplace climate.
    \item Observed sorting is then not a pure preference measure.
    \item Beliefs designs should name whether they measure expected treatment, shift beliefs, or observe intended choices.
\end{itemize}
\end{frame}

\begin{frame}{Later adjustment object}
\[
V_{it}(o)
=
\max
\left\{
u_{it}(o)+\beta V_{i,t+1}(o),
\max_{o' \neq o}\left[u_{it}(o')-K_{i,t}(o,o')+\beta V_{i,t+1}(o')\right]
\right\}
\]
\begin{itemize}
    \item \(K_{i,t}\): retraining, lost seniority, networks, credentials, relocation, flexibility constraints.
    \item Later adjustment is a separate channel from initial field choice.
\end{itemize}
\end{frame}

\begin{frame}{Later adjustment and inertia}
\begin{itemize}
    \item First placements can define training, networks, tasks, schedules, and promotion gates.
    \item Early-career shocks can alter long-run labor and family trajectories.
    \item High-skill career ladders often reward continuous hours and early promotion races.
    \item Direct barriers can make movement out of early tracks costly even when workers want to switch.
\end{itemize}
\end{frame}

\begin{frame}{Early-career path dependence}
\small
\begin{tabular}{p{0.25\linewidth} p{0.55\linewidth}}
\toprule
Object & Evidence to look for \\
\midrule
First job or specialty & quasi-random assignment, placement shock, cohort rules \\
Career ladder & wage growth, promotion, task assignment, retention \\
Flexibility constraint & hours, schedule, family timing, job switches \\
Persistent mismatch & later occupation, earnings, authority, job quality \\
\bottomrule
\end{tabular}
\end{frame}

\begin{frame}{Channels can interact}
\begin{itemize}
    \item Anticipated future constraints can change initial major or occupation choice.
    \item Early choices build specific skills, credentials, networks, and seniority.
    \item Family timing can raise the cost of retraining or switching later.
    \item Firm promotion systems can amplify small early differences into large career gaps.
\end{itemize}
\end{frame}

\begin{frame}{Short bridge to Week 5 norms}
\begin{itemize}
    \item Norms shape aspirations, acceptable jobs, self-promotion, mobility, and perceived feasibility.
    \item Domestic norms can affect expected care responsibilities before career entry.
    \item Workplace norms can affect whether technical, competitive, or authority-bearing roles feel accessible.
    \item Week 3 flags the bridge; Week 5 studies norms and institutions directly.
\end{itemize}
\end{frame}

\begin{frame}{Methods map}
\scriptsize
\setlength{\tabcolsep}{3pt}
\begin{tabular}{p{0.22\linewidth} p{0.23\linewidth} p{0.28\linewidth}}
\toprule
Design & Identifying variation & Observed margin \\
\midrule
Role-model intervention & classroom exposure & major or field choice \\
Teacher-bias design & teacher signals / grading environment & course, track, field \\
Expectations survey & measured or shifted beliefs & intended or realized choice \\
Early-career quasi-experiment & placement, shock, assignment & persistence, switching, wage growth \\
Long panel & lifecycle histories & timing, path dependence \\
Decomposition & accounting variation & shares across occupations or fields \\
\bottomrule
\end{tabular}
\end{frame}

\begin{frame}{Week 3 lab}
\begin{itemize}
    \item Reproduce: synthetic role-model intervention inspired by Porter and Serra.
    \item Diagnose: treatment, outcome, identifying variation, margin, and limitation.
    \item Transfer: synthetic competitiveness and track-choice exercise inspired by Buser, Niederle, and Oosterbeek.
    \item Optional frontier prompt: anticipated discrimination or early-career setbacks.
\end{itemize}
\end{frame}

\begin{frame}{Research frontier}
\begin{itemize}
    \item Which early-choice interventions produce durable labor-market outcomes?
    \item How do role models change information, belonging, expected returns, and discrimination beliefs?
    \item Which early-career shocks reveal true adjustment frictions?
    \item What data connect education, first job, firm assignment, promotion, and family timing?
\end{itemize}
\end{frame}

\begin{frame}{Bridge to Week 4}
\begin{itemize}
    \item Week 3 explains who enters which fields and early career ladders.
    \item Week 4 moves inside firms: hiring, evaluation, pay-setting, promotion, authority, and retention.
    \item The next question is how firms transform worker sorting into pay, authority, and career outcomes.
\end{itemize}
\end{frame}

\begin{frame}{Takeaways}
\begin{itemize}
    \item Gendered sorting has early-choice and later-adjustment channels.
    \item Observed sorting is not a pure preference or pure constraint statistic.
    \item Credible evidence names the variation, the observed margin, and the mechanism.
    \item Week 3 links household expectations to firm-side career dynamics and the later norms lecture.
\end{itemize}
\end{frame}

\end{document}
